\documentclass[11pt]{article}
\usepackage{geometry}\geometry{margin=1in}
\usepackage{setspace}\onehalfspacing
\usepackage{amsmath,amssymb}
\usepackage{xcolor}
\usepackage{hyperref}
\hypersetup{colorlinks=true, linkcolor=blue, citecolor=blue, urlcolor=blue}
\title{Response to Referee Comments\\
\large \emph{Sheltered Bidding: The Within-Auction Cost of SME Set-Asides}}
\author{Darcio Genicolo-Martins (INSPER)}
\date{April 2026}
\begin{document}
\maketitle

\noindent We are grateful to the referee for the detailed and constructive evaluation. The report identified eleven major concerns and five minor concerns. Each is addressed below with the specific revision and the section/page reference in the revised manuscript. Where the original specification is preserved, we provide the rationale for not making the requested change.

\section*{Major concerns}

\paragraph{1. Identification of the March 2018 cutoff (event study + placebo).}
We have added a six-semester event-study figure (Figure~F.2) and a placebo treatment-date table (Table~F.31) to the Reduced-Form Benchmark appendix. The placebo cutoffs at September 2017, March 2017, and June 2017 deliver price coefficients of $-0.013$ (insignificant), $-0.030$ (significant at 5\%), and $-0.034$ (significant at 1\%). The placebo evidence does not establish absence of pre-period correlation: low-frequency variation in the treated group is detectable at fake cutoffs near the real one. The relevant test is therefore magnitude separation: the actual March 2018 coefficient of $-0.108$ to $-0.142$ is three to eight times larger in absolute value than any placebo coefficient and is significant at the 1 percent level in every specification. The structural calibration is anchored to a treatment effect of demonstrably larger magnitude than any pre-period drift; the residual concern that some of the real-cutoff coefficient reflects gradual anticipation rather than discrete treatment is a bound, not a refutation, on the structural input.

\paragraph{2. Never-treated control group comparability.}
The placebo evidence in Table~F.31 disciplines the control set under the magnitude-separation reading: fake cutoffs produce coefficients several times smaller than the real one. The semester-binned event study (Figure~F.2) corroborates the post-cutoff shift in direction and magnitude. The 76 never-treated product groups remain the natural control set because (i) they are subject to the same federal SME-only default that already applied to all eligible items in those groups, and (ii) Group~65 is the only group with a regime change at the March 2018 cutoff. We have added a sentence in Appendix~F clarifying these points.

\paragraph{3. Invariance assumptions and KS rejection.}
We have added a paragraph in Section~3 (model) explicitly addressing this concern: \emph{the KS rejection of strict invariance is the empirical foundation for adopting equilibrium selection (Assumption~3) as the main specification rather than strict invariance}. A model that imposed strict invariance and then ran the BNE counterfactual on data that visibly violate it would be misspecified by construction. Section~7.1 has been reframed accordingly: strict invariance is the limiting bound against which the main specification is reported, not a competing baseline.

\paragraph{4. IPV restriction tested.}
The cross-modality consistency check in Section~5 (Table~B.6) and the new affiliation sensitivity exercise in Section~7 (Table~D.23, Gaussian copula with $\rho_c \in \{0, 0.1, 0.2, 0.3\}$) jointly bound the IPV restriction. The within-auction share moves by less than 5 percentage points and the total simulated effect by less than 10 percent under affiliation up to $\rho_c = 0.3$. Higher affiliation levels are inconsistent with the cross-modality test (which would reject under strong common-values contamination), so the relevant range for IPV-versus-affiliation sensitivity is bounded.

\paragraph{5. DiD-structural gap (2--3$\times$).}
We have added Table~C.15 (\emph{Bridge from DiD coefficient to structural counterfactual}) decomposing the gap into four mechanical sources: (a) sample restriction, (b) UH cleaning, (c) functional form (second-order statistic vs.\ linear log-mean), (d) conditioning set (counterfactual vs.\ realized entry). The four sources sum to the structural object up to a small numerical residual ($+0.004/+0.003$), demonstrating that the gap is mechanically reproducible and predictable in direction (structural $>$ DiD) under correct specification, not a sign of misspecification. Section~6.5 has been expanded to develop this point explicitly.

\paragraph{6. Methodological contribution overstated.}
We have softened the claim in the introduction (Section~1) and Section~3. The contribution is now positioned as the \emph{combination} of a clean total-exclusion legal trigger, Preg\~ao drop-out point-identification, Krasnokutskaya UH correction, observed-equilibrium-entry treatment, and Saez-Stantcheva welfare-weight identity. The combination delivers a within-auction structural decomposition that the partial-preference structural literature identifies only jointly with cross-equation constraints, not from a single legal trigger. We do not claim methodological invention of any individual component.

\paragraph{7. Welfare arithmetic and $\lambda$ sensitivity.}
We have added Table~E.27 (\emph{Welfare ranking stability across $\lambda$}) explicitly reporting the V0 vs V3 ranking at each $\lambda \in \{0.15, 0.20, 0.30, 0.40, 0.45\}$. The non-pharma ranking is V3 $>$ V0 across the full $\lambda$ range under both specifications. The pharma ranking is V3 $>$ V0 under the main specification and V0 $>$ V3 under strict invariance, both stable across the full $\lambda$ range. The ranking is therefore not driven by the choice of $\lambda$.

\paragraph{8. Sample selection and attrition.}
We have added Table~A.1 (\emph{Sample construction: from BEC raw extract to structural sample}) reporting stage-by-stage attrition with separate Pre/Post sub-sample reporting. Attrition is dominated by the Group-65 and Preg\~ao-only restrictions; the bid-level filter and $n \ge 2$ filter together drop 6.4 percent of remaining observations. There is no jump in attrition at the March 2018 cutoff; monthly auction counts move smoothly across the window.

\paragraph{9. Entry cost calibration ($\kappa^k$).}
The R\$0.11--R\$2.46 figures are per-auction zero-profit margins at the entry margin, not fixed costs of preparing a documentation packet. Section~5 now includes a footnote with a back-of-envelope external validation: BEC participation requires one fiscal certificate (CND, $\approx$R\$10 amortized over many auctions), one product registration form (free), and one power-of-attorney filing ($\approx$R\$5). For a supplier participating in 100 auctions per quarter, the marginal opportunity cost per auction is on the order of R\$0.10--R\$0.20, consistent with the calibrated $\kappa^{\text{SME}} =$ R\$0.11 in pharma.

\paragraph{10. Pharmaceutical bifurcation (diagnostic, not fragility).}
Section~8.1 now positions the bifurcation explicitly as the central policy claim: \emph{the paper does not assert universal preference-rule dominance; it identifies the conditions (thin pool, heterogeneous good) under which the standard ranking reverses, and documents that those conditions are met in only one of the two Group-65 sub-classes studied}. Of the eight (class~$\times$~type~$\times$~period) cells in the structural sample, the policy reversal between V0 and V3 obtains in exactly the two pharma cells under strict invariance and in zero cells under the main specification. This is a sharp policy partition, not a fragility we hedge against.

\paragraph{11. V3 under FPSB Maskin-Riley equilibrium.}
We have added Table~D.24 (\emph{V3 under Vickrey-equivalent vs.\ FPSB (Maskin-Riley) equilibrium}) calibrating the FPSB equivalent at $k=10\%$. The SME bid-function shift is approximately $-0.05$ and the non-SME shift approximately $+0.02$; the net effect on the second-order statistic is 0.22--0.27 percent of $p^{\mathrm{ref}}$, well within Monte Carlo noise. The V3 vs.\ V0 policy ranking is preserved in both classes. The conditional welfare ranking is robust to the FPSB adjustment.

\section*{Minor concerns}

\paragraph{1. Table~F.28 (price-to-reference) scale.}
Footnote added in Appendix~F clarifying that the price-to-reference outcome is on a level (not log) scale, and that the smaller magnitudes at 12-month and 18-month windows reflect the bounded support of the outcome and the absorption of the cap-distance dynamics by item fixed effects.

\paragraph{2. SME proxy audit (Table~A.4).}
Caption rewritten to clarify that the union proxy is preferred for type-completeness coverage (the BNE counterfactual requires assignment of every bidder to a type), not for MAE minimization. The trade-off is documented.

\paragraph{3. Entry count consistency.}
Verified: all entry counts in Section~3, Table~C.12, and Table~D.22 are aligned. The ``$n$ auctions (pooled)'' labeling introduced in earlier revisions is now consistent across all consumers.

\paragraph{4. ``Extensions not attempted'' (Section~9.5).}
Rewritten as constructive: three natural follow-on extensions (quantity response, dynamic capacity, GE SME surplus) are listed with explicit bounds on each. None changes the organizing lesson of the paper.

\paragraph{5. External validity language.}
``Suggestive beyond the Group-65 split'' replaced with ``consistent with the Group-65 split''. Generalization explicitly limited to ``structurally similar Brazilian procurement settings''.

\section*{Summary of revisions}

\noindent The revised manuscript is 77 pages. New material added in this revision: 4 new tables (sample flow, DiD-structural bridge, welfare ranking by $\lambda$, affiliation sensitivity, Maskin-Riley FPSB calibration; 5 in total counting the integrated event-study/placebo), 1 new figure (event-study around March~2018 cutoff), and 6 new defensive paragraphs distributed across Sections~3, 5, 6, 7, 8 and Appendix~F. The headline numbers (within-auction share 74.5\%/73.3\% in non-pharma/pharma; 50--60\% larger latent shock under fixed pool; 28.7\%/47.0\% welfare cost at $\lambda=0.30$; conditional welfare ranking V3 $>$ V0 in non-pharma robustly and in pharma under main specification only) are unchanged. The structural decomposition, the welfare arithmetic, and the conditional policy claim are unchanged in substance; the revision strengthens the empirical defense of each.

\end{document}
